\chapter[chap:phase3]{Phase 3: Branches und Merge}

\section{Ziel}

Parallele Arbeit ist lokal möglich. Branches lassen sich anlegen, wechseln und zusammenführen.

\section{Umfang}

\startitemize
\item Befehle \type{branch}, \type{checkout} und \type{log}.
\item Drei-Wege-Merge über die Merge-Basis, den besten gemeinsamen Vorfahren der beiden Seiten.
Die Funktion \type{merge_bases} liefert alle Merge-Basen zweier Commits.
Nach einem Criss-Cross-Merge, bei dem sich zwei Seiten gegenseitig einmischen,
sind es mehrere, und der Merge muss sie zusammenführen.
Die Funktion liest kein Commit-Datum, sondern geht die ganze Vorfahrenmenge beider Seiten durch.
\item Konflikte als Daten im Commit, nicht als Markierung in der Arbeitskopie.
Kapitel 2 beschreibt das Vorbild Jujutsu.
Ein Konflikt ist eine Summe von Trees: Der erste Tree wird genommen, und die Änderung von jedem
entfernten Tree zum hinzugefügten danach wird angewandt.
Gleiche hinzugefügte und entfernte Terme heben sich auf.
Deshalb verschachtelt sich ein Konflikt nie: Ein zweiter Branch, der auf einen offenen Konflikt
trifft, verlängert die Summe.
Bleibt ein Tree übrig, ist der Zustand ein gewöhnlicher Commit mit unverändertem Format und
unveränderter ID.
Alle Pfade eines Commits teilen die Zahl der Seiten; ein kleinerer Konflikt wird mit
aufhebenden Paaren aufgefüllt.
Der Text mit Markierungen entsteht aus einer einzigen Funktion der Summe.
Er soll nur auf der Platte und im \type{diff} erscheinen; der Snapshot soll Bytes vergleichen und
keine Markierungen lesen.
Noch ruft ihn nichts auf.
Die Regel, dass zwei Seiten mit gleicher Änderung sofort als gelöst gelten, gehört nicht zur
Algebra, weil sie beim erneuten Zusammenführen Information verliert.
Noch schreibt kein Befehl einen solchen Commit.
\item Berechnung des Merges ist Risikostufe 2 nach ADR-004.
\item Eigentümer der Arbeitskopie (Format 5, ADR-010): Ein Snapshot speichert nur die
Änderungen des Akteurs, dem die Arbeitskopie gehört; ein anderer übernimmt sie ausdrücklich
mit \type{--take-over} (Risikostufe 3). Umgesetzt; Kapitel 4 beschreibt die Regel,
der Anhang die Befehle. Die zweite Arbeitskopie baut darauf auf.
\stopitemize

\section{Offene Fragen}

\startitemize
\item Rebase: Ja oder nein in Phase 3?
\item Verhalten bei einem Branch, der in einem anderen Worktree ausgecheckt ist.
\item Abgeleitete Daten, etwa Indizes für \type{log} oder später Blame.
ADR-017 legt einen Cache unter \type{.stemma/derived/} fest, der nie zum
Objektmodell gehört. Er wird bei der ersten Abfrage aus den Werten der Eltern gefüllt, ein
Löschen kostet nur Zeit, und ein Purge entfernt ihn. Offen bleiben die erste Art, die
gespeichert wird, und die Bereinigung alter Einträge.
Ein Kandidat ist eine Generationsnummer, die die Suche nach der Merge-Basis früh beenden
würde; sie kommt erst, wenn eine Messung sie verlangt.
\stopitemize
